% Заключение. Версия 2 --- после получения данных Главы 4.
% Объем 2 страницы.
% Раздел «Направления дальнейших исследований» --- связным текстом,
% минимум 10 строк (правки друга-ревьюера).

\section*{Заключение}
\addcontentsline{toc}{section}{Заключение}

В работе проведена количественная оценка влияния выбора
коммуникационной топологии, механизма ее адаптации и~позиции
человека на~эффективность решения задач разных классов MAS LLM.
Цель достигнута~--- получены количественные характеристики для
пяти статических топологий, адаптивной мета-топологии в~двух
режимах маршрутизации, пяти ролей человека и~их~адаптивной
комбинации; выполнено 4230~завершенных запусков на~четырех корпусах
задач с~суммарной стоимостью~\$105{,}35.

RQ1 (от англ. Research Question) подтвержден~--- универсальной
статики-доминатора нет, победители на~четырех корпусах различаются
(Chain~--- HumanEval, Star~--- GSM8K, Debate~--- CommonGen и~DABench).
RQ3 подтвержден~--- HITL дает прирост $+0{,}033$ при~неоднородности
по~ячейкам.

RQ2 раскладывается на~центральный вывод~--- зависящее от~сложности
задачи Парето-преимущество адаптивной мета-топологии на~нетривиальных
корпусах с~поведением наименьшего риска (Утверждение~1)~---
и~поддерживающее эмпирическое наблюдение об~инверсии Парето-победителя
между развернутыми конфигурациями работника при~переходе с~gpt-oss-120b
на~Qwen-3-235B (Утверждение~2; см.~\ref{sec:analysis:rq}). RQ4
опровергнут с~интерпретацией <<универсальный отрицательный результат>>
(Утверждение~3)~--- адаптивный маршрутизатор роли повсеместно
не~дает положительного прироста поверх адаптивной мета-топологии;
эффект не~реализационный, а~в~самой гипотезе двухуровневой адаптации.

Ключевой результат работы~--- количественная демонстрация ценности
адаптивного дизайна на~нетривиальных задачах. На~корпусах, где
единичная статическая топология имеет реальные ограничения
(CommonGen, DABench), адаптивная мета-топология обеспечивает
Парето-выигрыш и~поведение наименьшего риска по~трем независимым
осям (наихудший случай качества, дисперсия качества, разброс
стоимости), сохраняя при~этом структурные свойства дизайна~--- единый
конфиг развертывания, интеграцию HITL-сигнала и~фазовое переключение
под-топологий. Результат получен при~контролируемых условиях
одного семейства работника и~устойчив к~основным сценариям
конфаунда. Поддерживающее наблюдение по~инверсии Парето-победителя
мотивирует практику кросс-конфигурационной валидации настроек MAS
перед~переносом на~новую развернутую конфигурацию работника;
механизм инверсии (семейство весов, глубина рассуждений, поддержка
структурированного вывода) в~работе не~устанавливается. Универсальный
отрицательный результат по~RQ4 дает конкретное проектное
ограничение~--- адаптивность на~двух уровнях одновременно (топология
и~роль) избыточна, достаточно одного уровня (топологии).

Ограничения~--- $n=45$ запусков на~ячейку $(\text{корпус},\text{режим})$
в~E3 и~E4 (на~этом объеме устойчиво подтверждены три эффекта);
кросс-конфигурационная валидация на~двух конфигурациях рабочего агента;
наблюдаемая инверсия Парето-победителя смешана сразу с~несколькими
переменными (семейство весов, глубина рассуждений, поддержка
структурированного вывода) и~в~работе не~разделяется; DABench
проявляет специфичное для~модели ограничение на~gpt-oss-120b
(см.~\ref{sec:exp:conf}); человек моделируется языковой моделью;
корпус не~покрывает ряд классов (диалог, перевод, обобщение длинных
текстов); судья~--- единственная модель из~единственного семейства.

\subsection*{Направления дальнейших исследований}

Прямое продолжение работы образует четыре взаимодополняющих направления.
Первое~--- расширение кросс-семейной валидации на~три и~более семейства
весов рабочего агента (Anthropic Claude, Google Gemini, Meta LLaMA-4)
с~приведением параметра глубины рассуждений к~единой шкале, что устранит
смешение эффектов семейства и~режима работы агента и~позволит проверить,
является ли наблюдаемая кросс-семейная неустойчивость политик
маршрутизации универсальным свойством адаптивных систем или специфичной
для~конкретной пары семейств. Второе~--- валидация симулятора человека
на~реальных участниках по~схеме планируемого эксперимента~E5
с~использованием шкалы~NASA-TLX (от англ. NASA Task Load Index~---
индекс когнитивной нагрузки NASA), латино-квадратного дизайна
на~восьми-двенадцати участниках и~отобранного подмножества
из~двенадцати задач, что позволит численно оценить корреляцию между
симулированной и~фактической оценкой когнитивной нагрузки.
Третье~--- разработка семейно-осознанного маршрутизатора, в~котором
политика выбора топологии и~роли явно учитывает семейство весов
рабочего агента, например, через предварительное дообучение
маршрутизатора на~трассах целевой модели или через переключение между
семейно-специфичными политиками. Четвертое~--- расширение корпуса
задач на~непокрытые классы (диалог, перевод, обобщение длинных
документов, мультимодальные задачи) и~проверка обобщения результатов
воронки на~новые классы.
